\section{历史发展}\label{sec:history}

\subsection{早期历史（1742--1900）}

1742年6月7日，哥德巴赫在给欧拉的信中提出了以下猜想：

\begin{quote}
    每个大于2的偶数都可以写成两个素数之和。
\end{quote}

欧拉在回信中表示他认为这个猜想是正确的，但无法给出证明。这一问题随后成为数论中最著名的未解之谜。（注：原信按当时把 $1$ 视为素数的习惯，表述为“每个大于 $2$ 的整数都是三个素数之和”；上引形式为欧拉改写后的现代表述。）

\subsection{20世纪初期（1900--1930）}

20世纪初，随着解析数论的发展，数学家们开始使用新的工具研究哥德巴赫猜想。

\begin{itemize}
    \item \textbf{哈代与李特尔伍德}（1923）：引入圆法，证明了在广义黎曼假设成立的条件下，几乎所有的偶数都满足哥德巴赫猜想\cite{hardy1923some}。
    \item \textbf{维诺格拉多夫}（1937）：使用圆法证明了充分大的奇数可以表示为三个素数之和，即弱哥德巴赫猜想对充分大的奇数成立\cite{vinogradov1937representation}。
\end{itemize}

\subsection{中期进展（1930--1970）}

这一时期见证了筛法的重大发展：

\begin{itemize}
    \item \textbf{布伦}（1920）：引入布伦筛法，证明了每个充分大的偶数都可以表示为两个殆素数之和\cite{brun1920sieve}。
    \item \textbf{塞尔伯格}（1947）：发展了塞尔伯格筛法，改进了布伦的结果\cite{selberg1949sieve}。
    \item \textbf{陈景润}（1966）：证明了每个充分大的偶数都可以表示为一个素数和一个不超过两个素数乘积之和\cite{chen1966representation}。
\end{itemize}

\subsection{现代发展（1970--至今）}

\begin{itemize}
    \item \textbf{里奇斯坦}（2001）：使用计算机验证了哥德巴赫猜想对所有小于 $4 \times 10^{14}$ 的偶数成立。
    \item \textbf{奥利维拉·e·席尔瓦、赫尔佐格与帕尔迪}（2014）：将验证范围扩展到 $4 \times 10^{18}$\cite{oliveira2014computational}。
    \item \textbf{黑尔弗戈特}（2013）：完全证明了弱哥德巴赫猜想——每个大于 $5$ 的奇数都是三个素数之和\cite{helfgott2013ternary}（此前数值验证与解析估计已覆盖到约 $10^{30}$ 以上的范围）。
\end{itemize}

\subsection{重要人物}

\begin{table}[htbp]
    \centering
    \caption{哥德巴赫猜想研究的重要贡献者}
    \label{tab:contributors}
    \begin{tabular}{@{}lll@{}}
        \toprule
        \textbf{数学家} & \textbf{时间} & \textbf{主要贡献} \\
        \midrule
        哥德巴赫 & 1742 & 提出猜想 \\
        欧拉 & 1742 & 首次系统研究 \\
        哈代 \& 李特尔伍德 & 1923 & 圆法 \\
        维诺格拉多夫 & 1937 & 三素数定理 \\
        布伦 & 1920 & 布伦筛法 \\
        陈景润 & 1966 & 陈氏定理 \\
        \bottomrule
    \end{tabular}
\end{table}
